\documentclass{article}
\usepackage[margin=25mm]{geometry}
\usepackage{amsmath,amssymb}
\usepackage{xurl}
\usepackage[colorlinks=true,linkcolor=blue,citecolor=blue,urlcolor=blue]{hyperref}
% Shared commands for notes in docs/paper-gaps/.

\DeclareUnicodeCharacter{2080}{\ifmmode{}_{0}\else\textsubscript{0}\fi}
\DeclareUnicodeCharacter{2081}{\ifmmode{}_{1}\else\textsubscript{1}\fi}
\DeclareUnicodeCharacter{2082}{\ifmmode{}_{2}\else\textsubscript{2}\fi}
\DeclareUnicodeCharacter{2099}{\ifmmode{}_{n}\else\textsubscript{n}\fi}

\newcommand{\Lean}{\textsc{Lean}}
\ExplSyntaxOn
\NewDocumentCommand{\leanid}{m}{
  \ifmmode
    \textsf{\detokenize{#1}}
  \else
    \group_begin:
    \urlstyle{sf}
    \tl_set:Nn \l_tmpa_tl {#1}
    \tl_replace_all:Nnn \l_tmpa_tl {\_} {_}
    \exp_args:NV \path \l_tmpa_tl
    \group_end:
  \fi
}
\ExplSyntaxOff
\providecommand{\tr}{\operatorname{tr}}
\newcommand{\tnleanIssueBase}{https://github.com/LionSR/TNLean/issues/}
\newcommand{\tnleanPrBase}{https://github.com/LionSR/TNLean/pull/}
\newcommand{\tnleanDocBase}{https://sirui-lu.com/TNLean/docs/}
\newcommand{\ghissue}[1]{\href{\tnleanIssueBase#1}{GitHub issue \##1}}
\newcommand{\ghpr}[1]{\href{\tnleanPrBase#1}{PR \##1}}
\newcommand{\doclink}[1]{\href{\tnleanDocBase#1.html}{\path{#1}}}

% Dirac notation and the identity operator, as in blueprint/src/macros/common.tex.
% \providecommand keeps note-local definitions authoritative.
\usepackage{dsfont}
\providecommand{\ket}[1]{|#1\rangle}
\providecommand{\bra}[1]{\langle #1|}
\providecommand{\braket}[2]{\langle #1 | #2 \rangle}
\providecommand{\ketbra}[2]{|#1\rangle\!\langle #2|}
\providecommand{\Id}{\mathds{1}}
\providecommand{\one}{\mathds{1}}
\providecommand{\C}{\mathbb{C}}
\providecommand{\MN}[1]{M_{#1}(\C)}
\providecommand{\MD}{M_D(\C)}

% Verdict marker: \gapnote{<kind>}{<status>}. Kind is the policy
% classification (clarification, local-correction, scope-restriction,
% unfaithful, false-source, open-gap); status is open, wip (elimination
% underway), resolved, or historical. Machine-read by the site index;
% prints a verdict line.
\newcommand{\gapnote}[2]{\par\noindent\textbf{Verdict:} #1 (#2).\par}

\providecommand{\gapnote}[2]{\par\noindent\textbf{Verdict:} #1 (#2).\par}
\title{The First Uniform Energy Estimate for ALVV Soft Truncation}
\author{}
\date{October 29, 2026}
\begin{document}
\maketitle
\gapnote{false-source}{open}

\section{The estimate and its quantifiers}

Arad, Landau, Vazirani and Vidick~\cite{Arad2016RigorousRG},
\href{https://arxiv.org/abs/1602.08828v2}{arXiv:1602.08828v2}, define
\begin{align}
    h_{m,t}(x)=t\sum_{k=1}^m\frac{(1-e^{-x/t})^k}{k},
    \qquad t\ge m\ge2,
    \label{eq:alvv_soft_scalar}
\end{align}
where $m$ is an integer. Let $H=H_J+H_{\overline J}$ be a chain
Hamiltonian with positive nearest-neighbour contractions. If $\varepsilon_J$
is the ground energy of the internal sum $H_J$, choose
$\varepsilon_J-10\le\varepsilon'_J\le\varepsilon_J$ and set
\begin{align}
    \widetilde H
      &=\varepsilon'_J I+h_{m,t}(H_J-\varepsilon'_J I)+H_{\overline J},&
    \varepsilon&=\varepsilon_{\overline J}+\varepsilon'_J.
    \label{eq:alvv_soft_physical}
\end{align}
The first energy estimate in Lemma~\path{lem:exp-approx},
equation~\path{eq:exp-approx-0}, asserts, for a normalized eigenvector
$H\psi=\lambda\psi$,
\begin{align}
    0\le\lambda-\langle\psi,\widetilde H\psi\rangle
       \le O\!\left(\frac{(\lambda-\varepsilon)^m}{m t^{m-1}}\right).
    \label{eq:alvv_soft_printed}
\end{align}
The source locations are \path{local_alg_arxiv_v2.tex}, lines 1154--1157
(scalar function), 1175--1185 (positive terms and regional convention),
1195--1201 (estimate), and 1216--1227 (its proof).

The usual uniform interpretation of (\ref{eq:alvv_soft_printed}) is false.
The three-qubit family below has fixed local norms, local dimension and
spectral gap. For every fixed admissible pair $(m,t)$, the ratio of its
actual energy defect to the displayed scale is unbounded. Thus even a
constant depending on these fixed parameters cannot establish the estimate
uniformly over the Hamiltonian family. The separate condition on $t$ in the
subsequent gap assertion does not qualify this first displayed estimate;
the counterexample also applies to arbitrarily large fixed parameters.

The counterexample is an exact mathematical calculation, not a formalized
Lean proof. It concerns this first uniform estimate only. It does not assert
that every later spectral, low-energy, or efficient-approximation conclusion
is false.

\section{A nondegenerate three-qubit family}

Let $R=|1\rangle\langle1|$, $P=|0\rangle\langle0|$, and, for
$1/2<c<1$, put
\begin{align}
    v_c&=c|0\rangle+\sqrt{1-c^2}|1\rangle,&
    Q_c&=|v_c\rangle\langle v_c|,\notag\\
    h_{12}&=\frac{R_1+P_2}{2},&
    h_{23}&=\frac{Q_{c,2}+R_3}{2},& H_c&=h_{12}+h_{23}.
    \label{eq:alvv_soft_example}
\end{align}
Identities on the other sites are suppressed. Each local term has spectrum
$\{0,1/2,1\}$, norm one and minimum zero, and acts nontrivially on both
sites of its assigned bond. These are admissible positive contractions.

The middle-qubit operator $(P+Q_c)/2$ has eigenvalues $(1-c)/2$ and
$(1+c)/2$. Hence the unique normalized ground vector and its energy are
\begin{align}
    \psi_c&=|0\rangle\otimes
      \frac{|0\rangle-v_c}{\sqrt{2(1-c)}}\otimes|0\rangle,&
    \lambda_c&=\frac{1-c}{2}.
    \label{eq:alvv_soft_ground}
\end{align}
An outer-qubit excitation costs $1/2$, and the middle-qubit excitation costs
$c>1/2$. Thus the full spectral gap is exactly $1/2$ for every $c$.
Directly from (\ref{eq:alvv_soft_ground}),
\begin{align}
    \langle\psi_c,P_2\psi_c\rangle
      =\langle\psi_c,Q_{c,2}\psi_c\rangle=\lambda_c.
    \label{eq:alvv_soft_expectations}
\end{align}

Take the contiguous region $J=\{1,2\}$. The source convention gives
$H_J=h_{12}$ and $H_{\overline J}=h_{23}$: the crossing interaction is
included in the complementary remainder, not discarded. Both regional sums
have ground energy zero. Thus $\varepsilon'_J=0$ is admissible and
$\varepsilon=0$ in (\ref{eq:alvv_soft_physical}).

\section{The exact defect at every fixed truncation scale}

Differentiating (\ref{eq:alvv_soft_scalar}) and using $h_{m,t}(0)=0$ gives
\begin{align}
    x-h_{m,t}(x)
      =t\int_0^{x/t}(1-e^{-u})^m\,du,\qquad x\ge0.
    \label{eq:alvv_soft_defect_integral}
\end{align}
Consequently $a_{m,t}=1/2-h_{m,t}(1/2)$ is strictly positive for every
fixed $m\ge2$ and $t\ge m$. On the sector in which the first qubit is
$|0\rangle$, $h_{12}=P_2/2$. Functional calculus and
(\ref{eq:alvv_soft_expectations}) therefore give the exact identity
\begin{align}
    \lambda_c-\langle\psi_c,\widetilde H_c\psi_c\rangle
      =a_{m,t}\lambda_c.
    \label{eq:alvv_soft_exact_defect}
\end{align}
For the explicitly admissible choice $m=t=2$,
\begin{align}
    a_{2,2}=4e^{-1/4}-e^{-1/2}-\frac52>0,
    \label{eq:alvv_soft_small_scale}
\end{align}
where positivity follows from (\ref{eq:alvv_soft_defect_integral}).
The displayed source scale is then $\lambda_c^2/4$, and the ratio is
$4a_{2,2}/\lambda_c\to\infty$ as $c\uparrow1$.

More generally, for every fixed admissible $(m,t)$ the ratio is
\begin{align}
    \frac{a_{m,t}\lambda_c}{\lambda_c^m/(m t^{m-1})}
       =m t^{m-1}a_{m,t}\lambda_c^{1-m}\longrightarrow\infty.
    \label{eq:alvv_soft_divergence}
\end{align}
All parameters except $c$ are fixed in this limit. Each individual
Hamiltonian has a unique ground state and gap $1/2$; no degeneracy or
vanishing spectral gap is involved.

\section{The failed inference and the remaining comparison}

The scalar estimate gives the valid moment bound
\begin{align}
    0\le\langle\psi,(H-\widetilde H)\psi\rangle
      \le\frac{\langle\psi,A^m\psi\rangle}{m t^{m-1}},
    \qquad A=H_J-\varepsilon'_J I\ge0.
    \label{eq:alvv_soft_moment}
\end{align}
The source then uses $0\le A\le H-\varepsilon I$ to replace this
moment by $(\lambda-\varepsilon)^m$. Order of noncommuting positive
operators does not justify that replacement. In the present example,
$A=h_{12}\le H_c$, but
\begin{align}
    \langle\psi_c,A^2\psi_c\rangle=\lambda_c/4
      >\lambda_c^2=\langle\psi_c,H_c^2\psi_c\rangle,
    \label{eq:alvv_soft_failed_power_order}
\end{align}
because $0<\lambda_c<1/4$.

A corrected physical spectral comparison needs a separate argument.
One possible route is to retain moments of the actual regional operator
and derive their control from exponential moments. For example, scalar
functional calculus gives, for $s>0$,
\begin{align}
    \|(H-\widetilde H)\psi\|
      \le\frac{\|A^m\psi\|}{m t^{m-1}}
      \le\frac{(m/(es))^m}{m t^{m-1}}\|e^{sA}\psi\|.
    \label{eq:alvv_soft_possible_route}
\end{align}
A uniform exponential-moment estimate has now been proved for every
prescribed interval of a positive unit nearest-neighbour chain. After
subtracting the three attained regional minima, write $E_0$ for the
centered ground energy and $G_\eta$ for its spectral subspace below
$E_0+\eta$. For $\eta\ge0$ and
$\varepsilon_J-10\le\varepsilon'_J\le\varepsilon_J$, the bound is
\begin{align}
    \bigl\|e^{(H_J-\varepsilon'_J I)/8}|_{G_\eta}\bigr\|
       \le 5e^{(17+\eta)/8}.
\end{align}
The construction includes empty intervals and endpoints, and assumes
neither a spectral gap nor a ground-space dimension. For equal order and
scale, the same construction now gives the actual soft Hamiltonian
$\widetilde H_{m,m}\le H$ and
\begin{align}
    \bigl\|(H-\widetilde H_{m,m})|_{G_\eta}\bigr\|
        \le80e^{(17+\eta)/8-m/16}.
\end{align}
For $\alpha=(m/2)2^{-m}$ and $m/2\ge10$, the actual centered soft
ground energy can also be chosen before the bandwidth, with
$-\alpha\le E_s\le E_0$. If $m/2\ge17+\eta+\alpha$, its own low spectral
space $\widetilde G_\eta$ satisfies the corrected reverse residual bound
\begin{align}
    \bigl\|(H-\widetilde H_{m,m})|_{\widetilde G_\eta}\bigr\|
        \le80e^{(17+\eta+\alpha)/8-m/16}.
\end{align}
This follows from shifted soft/hard coercivity and an exponential moment on
the actual soft window. For $\eta,w,\delta>0$, set
\begin{align*}
    \rho&=\min\left\{1,\frac w2,
                         \frac{\delta w^2}{2(\eta+2)}\right\},&
    m&=\left\lceil35+2\eta+16\log(80/\rho)\right\rceil.
\end{align*}
This order is uniform in the chain, local dimension and Hamiltonian. Both
actual residuals are at most $\rho$, and $0\le E_0-E_s\le\rho$.
The soft space below $E_s+\eta+w$ is $\delta$-close to the original space
below $E_0+\eta$, and the original space below $E_0+\eta+w$ is
$\delta$-close to the soft space below $E_s+\eta$. These corrected
comparisons use the same original minima and the actual eigenbases. They
concern one interval, equal scale and order, and positive bandwidth and
enlargement. The equal-scale conclusions are instances of the bounded independent-scale
extension below. The printed gap prescription and the complete efficient
multi-region construction remain separate.


\subsection{A compatible independent-scale construction}

The exact identity
\begin{align*}
    h_{m,t}(x)=\frac tm h_{m,m}\left(\frac mt x\right)
\end{align*}
and the weighted scalar remainder imply, for $m\geq1$, $0<t\leq4m$ and
$x\geq0$,
\begin{align*}
    0\leq x-h_{m,t}(x)\leq16e^{x/8-t/16}.
\end{align*}
The same physical moment therefore proves the original and soft-window
residuals at independent scale, with decay $e^{-t/16}$, correction
$\alpha=(t/2)2^{-m}$, and actual hard comparator cutoff $t/2-10$.
No printed power estimate is used.

For $\eta,w,\delta>0$ and the same $\rho$ as above, choose
\begin{align*}
    m=\left\lceil12+\frac23\eta+
                     \frac{16}{3}\log(80/\rho)\right\rceil,
    \qquad t=3m.
\end{align*}
Then $\alpha\leq3/4$, the reverse cutoff condition holds, and both
residuals are at most $\rho$. The resulting actual soft Hamiltonian has
both directed window comparisons and actual operator tensors at any
prescribed norm accuracy. The heat times $a/t$, $0\leq a\leq m$, are
all at most $1/3$ and hence lie strictly within the corrected cluster
regime. The construction represents the exterior terms and scalar exactly,
and uses direct heat approximants for the regional term. Its explicit
bond budget is stated in the cluster-gap note linked below.

The order is selected before the physical chain and Hamiltonian; the
attained original minima precede the approximate interval energy, and the
actual soft minimum precedes the bandwidth. These conclusions concern one
prescribed interval. Simultaneous replacements on several regions, the
complete improving polynomial family and its algorithmic running time
are not supplied by this one-interval result.

This obstruction is distinct from the incorrect cluster-approximation
estimate, recorded in the separate note
\href{https://sirui-lu.com/TNLean/paper-gaps/alvv16_cluster_approximation.pdf}
{\emph{Support Size and Negative Local Terms in the ALVV Cluster Approximation}}.
The corrected scalar soft-truncation estimates are not refuted here, and no
failure of every subsequent conclusion follows from this counterexample.
The unrestricted soft-truncation conclusions remain open under
\href{https://github.com/LionSR/TNLean/issues/9847}{issue 9847} and the
\href{https://github.com/LionSR/TNLean/issues/9833}{ALVV tracker}.

\bibliographystyle{unsrt}
\bibliography{references}
\end{document}
